\begin{claim}{F.11}{Proved}{The rank-six cyclic centralizer and an exact common kernel}
Let $P=P\Gamma_{0,12}$ for the specified geometric
icosahedral action in F.7.  Fix $D=C_5\leq A_5$ and
$N=N_{A_5}(D)\cong D_{10}$, where $D_{10}$ has order ten.
The explicit configuration and marking below give
\[
 \begin{split}
 C_P(D)&=C_P(N)\cong\pi_1(X_5)=F_6,\\
 C_{\widehat P}(D)&=C_{\widehat P}(N)
                  \cong\widehat{\pi_1(X_5)}=\widehat F_6,
 \qquad
 X_5=\mathbf P^1\setminus\{0,\infty,\mu_5\}.
 \end{split}                                                   \tag{F.11.1}
\]
These are the actual discrete and completed inclusions
constructed below.  Of the twelve seven-point restrictions
in F.7.9, two are trivial on this subgroup; the other ten
have exactly the same kernel.  Their joint kernel, including
all twenty-four six-point restrictions, is the kernel of
the specified map
\[
 f:X_5\longrightarrow M_{0,7},\qquad
 t\longmapsto
 (1,\zeta^2,\zeta^4,t,t\zeta,t\zeta^4,\infty),
 \qquad \zeta=e^{2\pi i/5}.                                  \tag{F.11.2}
\]
The equalities of kernels hold for the discrete groups
and their full profinite completions.

At $b=2+2i$, the precisely based meridians below give
\[
          w=[m_{\zeta^2},\,m_\zeta m_{\zeta^3}m_0]
                    \in\ker(f_*),\qquad w\ne1.               \tag{F.11.3}
\]
Its procyclic closure is a full $\widehat\Z$ in this
common completed kernel.  After the specified basepoint
transport and inclusion into $\widehat P$, that
procyclic subgroup has trivial intersection with
$C_{\widehat P}(A_5)$, by F.10.  Thus the restrictions
again have a concrete common kernel, while the residual
centralizer in F.7 remains unresolved.
\end{claim}

\begin{proof}
\emph{The actual action, labels, and basepoint transport.}
Use the same five-letter generators and based coset
labels as F.7.19--57:
\[
 a=(1\ 2)(3\ 4),\qquad b=(1\ 3\ 5),\qquad d=ab,\qquad
                       D=\langle d\rangle .
\]
The indices $0,\ldots,11$ refer to the twelve cosets
of $D$, in that previously specified order.  Put
\[
 \theta=2\sin(2\pi/5),\quad
 \theta^4-5\theta^2+5=0,\quad 19/10<\theta<2,\qquad
 k=1-\theta^2,\qquad t_0=k\zeta^2 .
                                                               \tag{F.11.4}
\]
The exact marked configuration is
\[
 \begin{array}{c|ccccc}
 j&0&1&2&3&4\\ \hline
 \text{label of }\zeta^j&1&3&5&6&2\\
 \text{label of }t_0\zeta^j&4&7&9&11&8
 \end{array}
 \qquad 0\longleftrightarrow0,\quad
                      10\longleftrightarrow\infty .           \tag{F.11.5}
\]
Here are actual M\"obius transformations realizing the
named generators and the normalizer:
\[
 \begin{aligned}
 a(z)&=\frac{1-z}{1-z/k},&
 b(z)&=a(\zeta z),\\
 d(z)&=\zeta z,&
 \tau(z)&=\frac{t_0}{z},\\
 \tau&=(1\ 5)(3\ 4)\in A_5,&
 \tau d\tau^{-1}&=d^{-1}.
 \end{aligned}                                                 \tag{F.11.6}
\]
The finite certificate
\path{code/part_a/c5_centralizer_labels_f11.py}
checks these formulas in
$\mathbf Q(\theta,i)$, the $720$ images of the
twelve points under the sixty matrices, and all
$120$ products by the two generators.  The sixty
distinct matrices therefore form exactly the given
$A_5$, with precisely its stated coset action.
In particular this supplies an actual equivariant
identification of the marked action, not just an
isomorphism of abstract groups.  It is the explicit
M\"obius realization of the based covering choice
in F.7; changing that initial marking conjugates
all the inclusions and all the named forgetting
maps by the same marking isomorphism.

For every $t\in X_5$ replace $t_0$ in (F.11.5) by $t$.
This gives a map
\[
 j:X_5\longrightarrow M_{0,12},\qquad
                   t\longmapsto(0,\infty,\mu_5,t\mu_5)
                                                               \tag{F.11.7}
\]
with the individual labels in (F.11.5).
The functions $d(z)=\zeta z$ and $\tau_t(z)=t/z$
give a labelled $D_{10}$ action on this entire
family: $\tau_t$ exchanges the two poles and
sends the first orbit point with index $j$ to
the second with index $-j$, and conversely.
They satisfy the group relations on every fibre.
Monodromy therefore centralizes the actual
$D_{10}$ at $t_0$.
Equivalently, each of the ten corresponding
label permutations fixes the map to labelled
moduli (F.11.7), through these specified M\"obius
transformations, and hence fixes every induced
based loop.

The straight segment $\gamma$ from $t_0$ to
$b=2+2i$ lies in $X_5$: the exact field calculation
gives $\operatorname{Re}t_0>2$, so the whole
segment has real part at least two.
The element of $C_P(N)$ associated to (F.11.3)
is explicitly
$j_*(\gamma w\gamma^{-1})$.
All loop products here mean chronological
concatenation.  This also fixes the basepoint
transport for the seven- and six-point target maps.

\emph{The exact discrete centralizer and its full topology.}
The quotient by $d$ is the actual cover
$z\mapsto z^5$; its four punctures are
$0,\infty,1,t_0^5$.  An equivalent based group
description is
\[
 L=\pi_1\bigl(\mathbf P^1\setminus\{0,\infty,1,t_0^5\}\bigr)
       =F_3,\qquad
 U=\ker(\chi:L\longrightarrow C_5)=\pi_1(S_{0,12}),
                                                               \tag{F.11.8}
\]
where the positive meridians at $0,\infty$ have
$\chi$-values $1,-1$ and the other two have value
zero.  The inclusion and the quotient map in
(F.11.8) are induced by this covering.

We spell out why descent accounts for the entire
mapping-class centralizer.  Conjugation embeds
$L$ in $\operatorname{Aut}(U)$: an element of the
free group $L$ centralizing the nonabelian subgroup
$U$ is trivial.  Its image in $\operatorname{Out}(U)$
is the specified deck group $D$.
For an actual automorphism $\alpha$ representing
a pure mapping class centralizing $D$, conjugation
by $\alpha$ preserves this embedded $L$.
Indeed its effect on every deck outer class is
the identity, so its effect on the conjugation
automorphism of an $l\in L$ differs from that
automorphism by conjugation by an element of $U$.
The embedding of $L$ makes the resulting element
unique.  This defines an automorphism
$\alpha_L$ of $L$, restricting to $\alpha$ on $U$
and inducing the identity on $C_5$.

The two ramified cusp meridians in $U$ are fifth
powers of the corresponding meridians of $L$;
the other cusp meridians are their specified
unramified lifts.  Since $\alpha$ fixes each
upstairs peripheral class with exponent one,
the uniqueness of roots in a discrete free group
shows that $\alpha_L$ fixes each of the four
downstairs peripheral classes with exponent one.
The discrete Dehn--Nielsen--Baer identification
used in A.12 therefore gives a pure mapping class
on the four-punctured quotient.
Changing $\alpha$ by an inner automorphism of $U$
does not change that downstairs class.
The kernel of this descent consists of deck
classes, since an inner automorphism of $L$
restricts to precisely such a class on $U$.
A deck class fixing the upstairs labels is the
identity, because $D$ acts freely on each
five-cycle.  This proves injectivity of descent
on the pure centralizer.

Conversely, every pure class on the quotient
preserves $\chi$ and has five commuting lifts.
Choose the unique lift fixing the labelled
point $1$ of the first five-cycle; it then fixes
that entire cycle.  Identify
$P\Gamma_{0,4}$ with
$\pi_1(\mathbf P^1\setminus\{0,1,\infty\},t_0^5)$
by pushing its fourth point while the first
three are fixed, as in A.12.
The chosen lift fixes the second five-cycle
exactly when this loop lifts to a closed loop
under $t\mapsto t^5$.
Its image is consequently the kernel of winding
about zero modulo five, namely
$\pi_1(X_5,t_0)$.
The lift is exactly the family (F.11.7).
This proves
$j_*\pi_1(X_5)=C_P(D)$ and identifies its source
as the index-five subgroup $F_6$ of this specified
$F_2$.
Since its monodromy also centralizes $\tau$,
it proves $C_P(D)=C_P(N)$.

There is a particularly direct check of the full
induced topology.  Retain the eight labels
\[
              I_8=\{0,10\}\cup\{1,3,5,6,2\}\cup\{4\}.
                                                               \tag{F.11.9}
\]
After this forgetting map, (F.11.7) has the seven
fixed points $0,\infty,\mu_5$ and the single
moving point $t$.  Filling that moving point
is precisely the normalization map of A.12.
On its point-pushing subgroup the resulting
automorphism of $\pi_1(X_5)$ is inner conjugation
by the source loop.  After completion the exact
composite is therefore
\[
 \widehat{\pi_1(X_5)}
 \xrightarrow{\widehat j_*}\widehat P
 \longrightarrow\widehat{P\Gamma}_{0,8}
 \longrightarrow\operatorname{Aut}(\widehat{\pi_1(X_5)}),
               \qquad u\longmapsto\operatorname{Inn}(u).
                                                               \tag{F.11.10}
\]
The final extension is continuous on each finite
characteristic quotient.  Its displayed restriction
is injective because the free profinite group has
trivial centre.  Thus $\widehat j_*$ is injective
and its image is exactly the closure of the
discrete centralizer.  The prime-order centralizer
comparison already justified in F.7 identifies
that closure with $C_{\widehat P}(D)$.
Taking closures of $C_P(D)=C_P(N)$ gives both
inclusions needed for the second equality in
(F.11.1).

\emph{All thirty-six restrictions and the actual normalizer transports.}
For $D$ itself the two seven-sets are
\[
 \begin{split}
 I_A&=\{0,1,2,3,5,6,10\},\\
 I_B&=\{0,4,7,8,9,10,11\}.
 \end{split}                                                   \tag{F.11.11}
\]
Their maps to labelled moduli are constant:
the first configuration is already fixed, and
the second becomes fixed under $z\mapsto z/t$.
For the other ten sets choose
$I=\{1,2,4,5,7,8,10\}$ as representative.
In the ordered target labels
$(1,5,2,4,7,8,10)$ its map is exactly (F.11.2).
The following are the actual $D_{10}$ transports:
\[
\begin{array}{c|l|l}
 k& d^k I&d^k\tau I\\ \hline
 0&\{1,2,4,5,7,8,10\}&\{0,1,2,3,4,7,11\}\\
 1&\{1,3,4,6,7,9,10\}&\{0,1,3,5,7,8,9\}\\
 2&\{2,3,5,7,9,10,11\}&\{0,3,4,5,6,9,11\}\\
 3&\{1,5,6,8,9,10,11\}&\{0,2,5,6,7,8,11\}\\
 4&\{2,3,4,6,8,10,11\}&\{0,1,2,4,6,8,9\}.
\end{array}                                                     \tag{F.11.12}
\]
The label certificate enumerates the six order-five
subgroups, verifies this orbit, and records each
transporter as a five-letter permutation, a word
in $a,b$, and its permutation of the twelve labels.
For $n=d^k\tau^\epsilon$, the actual homeomorphism
in (F.11.6) gives the target isomorphism
$n_I:P\Gamma_{0,I}\to P\Gamma_{0,nI}$.
Naturality of forgetting gives
\[
        r_{nI}\circ\operatorname{conj}(n)
                           =n_I\circ r_I .
                                                               \tag{F.11.13}
\]
Conjugation by $n$ is the identity on the subgroup
(F.11.1), so these ten restrictions have equal
kernels there.  The same identity extends to
the full completions.

Each of the twenty-four six-sets is obtained by
removing one of the two poles from its specified
seven-set.  The certificate lists all twelve
pairs of these literal subset inclusions.
Forgetting that pole gives the exact factorization
of the corresponding six-point map through the
seven-point map.  Hence adding all twenty-four
constraints does not change the joint kernel.
Together with (F.11.11)--(F.11.13), this proves
the asserted joint-kernel equality for both
the discrete and completed groups.

\emph{Precisely based loops and exact monodromy.}
Put $r=31/1000$.  For
$q\in\{\zeta^2,\zeta,\zeta^3,0\}$,
define $m_q$ by the six successive straight
segments through
\[
 b,\quad q+r+ri,\quad q-r+ri,\quad
 q-r-ri,\quad q+r-ri,\quad q+r+ri,\quad b.
                                                               \tag{F.11.14}
\]
The square is counterclockwise, with its
whisker traversed in reverse at the end.
The certificate checks that every segment
avoids $0,\mu_5$, and that each square contains
exactly its named puncture.

Remove the last, fixed point $\infty$ from
(F.11.2).  This gives an actual ordered
configuration map to $\operatorname{Conf}_6(\mathbf C)$.
Appending $\infty$ and taking labelled M\"obius
classes recovers $f$.
Its loop map therefore factors through
$P_6\leq B_6$, with the injective inclusion
given by the ordered-to-unordered covering.
Multiply all six finite coordinates by
$1+(371/1000)i$, and project to their real parts.
The increasing base order of the six positions,
numbered $0,\ldots,5$, is
$(4,1,2,0,3,5)$.
Use the convention of F.9: a positive
$\sigma_j$ has its left strand above its
right strand just before the crossing,
and crossings are read chronologically.
Identifying the base chamber with the real
ordered configuration fixes the braid
generators up to its unique chamber isotopy.

The program
\path{code/part_a/c5_seven_kernel_exact.py}
uses fractions in the exact field
$\mathbf Q[\theta]/(\theta^4-5\theta^2+5)$,
with
\[
             \zeta=\frac{\theta^2-3+i\theta}{2}.
                                                               \tag{F.11.15}
\]
The polynomial is strictly increasing on
$(19/10,2)$ and has opposite endpoint signs.
Rational bisection isolates its specified
root; exact interval evaluation certifies
each nonzero field sign.
For every one of the $24$ source segments
and every one of the $15$ coordinate pairs,
the projected difference is affine in
the segment parameter.  Its endpoint signs
enumerate all crossings, and its single
possible root is computed exactly in the field.
Every crossing has nonzero imaginary
separation.  Distinct crossings have distinct
parameters, the strands are adjacent in
the current order, and the final ordering
agrees with direct endpoint evaluation.
The exact source-segment tests also exclude
every configuration collision.

The output contains $360$ pair checks,
$58$ simple crossings, and $2530$ certified
algebraic signs, with a rational isolating
interval obtained by $128$ bisections.
Thus the words describe the actual paths:
between crossings the ordering chambers
are convex; each crossing has exactly the
certified adjacent pair and sign.
This is the projection-isotopy argument
of Artin's Theorem~3 \cite{A-Artin1947}.
Writing signed integers for braid generators,
the words for $m_{\zeta^2}$ and for
$m_\zeta m_{\zeta^3}m_0$, after free reduction,
are respectively
\[
\begin{split}
 A={}&(-4,-3,-5,-1,-4,-2,-1,-1,2,4,1,5,3,4),\\
 B={}&(-4,-3,-5,-1,-4,-2,-3,-2,-1,-4,\\
     &\hspace{9mm}-1,-4,-3,-2,-3,4,-5,-1,3,4).
\end{split}                                                     \tag{F.11.16}
\]
The complete words for all four individual
meridians, of lengths $14,8,20,16$, and
every crossing are retained in the output.

Use the faithful Artin substitutions of F.9.8,
with the stated inverse substitution for
each negative generator.  Integer free-word
substitution and cancellation give
\[
       \operatorname{Artin}(AB)
           =\operatorname{Artin}(BA),\qquad
       \operatorname{Artin}(ABA^{-1}B^{-1})(X_i)=X_i
                    \quad(1\leq i\leq6).                       \tag{F.11.17}
\]
The output records the six common images
and the six single-letter final images.
Artin's faithfulness theorem, Theorem~14
and formulas (14)--(15) of the original
article \cite{A-Artin1947}, proves
$[A,B]=1$ in $B_6$.
The actual configuration factorization
therefore proves $f_*(w)=1$.

\emph{Nontriviality and the full procyclic subgroup.}
Fill all source punctures except
$\zeta^2,\zeta,\infty$.
This is an actual map to a twice-punctured
plane.  The loops $m_{\zeta^3}$ and $m_0$
become trivial; the other two loops have
winding vectors $(1,0)$ and $(0,1)$.
In the integral Heisenberg quotient with
product
\[
 (a,b,t)(a',b',t')
       =(a+a',b+b',t+t'+ab'),
\]
their images consequently have the form
$(1,0,s)$ and $(0,1,t)$.
Their commutator is exactly $(0,0,1)$,
for every $s,t\in\Z$.
Thus the image of $w$ is nontrivial.
For each prime $p$, reducing these integer
matrices modulo every $p^n$ extends this
homomorphism continuously to the full
source completion and then to
$\operatorname{UT}_3(\Z_p)$.
The central coordinate of $w^u$ is $u_p$.
All prime coordinates are therefore
detected, so
$u\mapsto w^u$ embeds $\widehat\Z$.

The already proved full embedding
(F.11.10) preserves this conclusion for
$j_*(\gamma w\gamma^{-1})$ in $\widehat P$.
Its whole procyclic closure lies in all
thirty-six kernels by continuity.
Finally this is a nontrivial discrete
pure mapping class, so the general
procyclic exclusion of F.10 applies and
gives the asserted trivial intersection
with $C_{\widehat P}(A_5)$.
No element of that residual centralizer
is produced, and no equality between it
and this joint kernel is asserted.
\end{proof}
